\documentclass[12pt]{article}

\setlength{\parindent}{0pt}
\setcounter{secnumdepth}{3}
\usepackage[margin=1in]{geometry}

% packages
\usepackage{amsmath, amsfonts, amsthm, amssymb}
\usepackage{graphicx, float}
\usepackage{mathtools}
\usepackage{titlesec}
\usepackage{interval}
\usepackage{hyperref}
\usepackage{siunitx}
\usepackage{titling}
\usepackage{vwcol}
\usepackage{setspace}
\usepackage{empheq}
\usepackage{cancel}
\usepackage{esdiff}
\usepackage{multicol}
\usepackage{mdframed}
\usepackage{esdiff}
\usepackage{tikzsymbols}
\usepackage{multicol}
\usepackage{tikz}
\usepackage{varwidth}
\usepackage{pgfplots}
\usepackage{fancyhdr}
\usepackage{enumitem}
\usepackage{tcolorbox}

\intervalconfig {
	soft open fences
}

% header
\pagestyle{fancy}
\fancyhf{}
\lhead{Ethan Fan}
\rhead{AP Calculus BC}
\cfoot{\thepage}

% section format
\titleformat{\section}
  {\large \bfseries}
  {}
  {0em}
  {}
\titlespacing*{\section}{0pt}{0.5em}{0.3em}

\titleformat{\subsection}[runin]
  {\bfseries}
  {}
  {0em}
  {}
\titlespacing*{\subsection}{0pt}{0pt}{0em}

\titleformat{\subsubsection}[runin]
  {\itshape}
  {(\roman{subsubsection})}
  {0em}
  {}

% commands
\newcommand{\alignedintertext}[1]{%
  \noalign{%
    \vskip\belowdisplayshortskip
    \vtop{\hsize=\linewidth#1\par
    \expandafter}%
    \expandafter\prevdepth\the\prevdepth
  }%
}

\newcommand*{\paren}[1]{\ensuremath\left(#1\right)}
\newcommand*{\limit}[2][x]{\ensuremath{\displaystyle\lim_{#1\to#2}}}
\newcommand*{\Limit}[3][x]{\ensuremath{\displaystyle\lim_{#1\to#2}\left[#3\right]}}
\newcommand*{\deriv}[1][x]{\ensuremath{\dfrac{\mathrm{d}}{\mathrm{d}#1}}}
\newcommand*{\Deriv}[2][x]{\ensuremath{\dfrac{\mathrm{d}}{\mathrm{d}#1}\left[#2\right]}}
\newcommand{\dydx}{\dfrac{dy}{dx}}
\newcommand{\dydt}{\dfrac{dy}{dt}}
\newcommand{\dxdt}{\dfrac{dx}{dt}}
\newcommand*{\iinteg}[2][x]{\ensuremath{\displaystyle\int #2\;\mathrm{d}#1}}
\newcommand*{\dinteg}[4][x]{\ensuremath{\displaystyle\int_{#2}^{#3}#4\;\mathrm{d}#1}}
\newcommand*{\abs}[1]{\ensuremath{\left|#1\right|}}
\newcommand*{\eps}{\varepsilon}
\newcommand*{\real}{\ensuremath\mathbb{R}}
\newcommand*{\integer}{\ensuremath\mathbb{Z}}
\newcommand*{\posint}{\ensuremath\mathbb{Z^+}}
\newcommand*{\cbrt}[1]{\ensuremath{\sqrt[3]{#1}}}
\newcommand*{\lheqzero}{\ensuremath{\underset{\text{L'H}}{\overset{\left[\frac00\right]}{=}}}}
\newcommand*{\lheqinfty}{\ensuremath{\underset{\text{L'H}}{\overset{\left[\frac{\infty}{\infty}\right]}{=}}}}
\newcommand{\tor}{\text{ or }}
\newcommand{\floor}[1]{\lfloor #1 \rfloor}

\newcommand{\vem}{\vspace{0.5em}}
\newcommand{\example}[1]{\section{Example #1}}
\newcommand{\problem}[1]{\section{Problem #1}}
\newcommand{\subproblem}[1]{\subsection{(#1)} \,}

\DeclareMathOperator{\dne}{DNE}
\DeclareMathOperator{\sgn}{sgn}

\newcommand{\definition}[1]{\begin{tcolorbox}[colback=red!5!white,colframe=red!75!black,parbox=false] #1 \end{tcolorbox}}

\newcommand{\theorem}[2]{\begin{tcolorbox}[title={#1},colback=blue!5!white,colframe=blue!75!black,parbox=false] #2 \end{tcolorbox}}

\allowdisplaybreaks
% \postdisplaypenalty=100000

% document
\begin{document}

\vspace*{0cm}

% opening
\begin{center}
    {\Large \bfseries Problem Set 3} \\[0.5cm]
    {\large Upper and Lower Sums/Convergence} \\[0.35cm]
    {\normalsize September 9, 2026}
\end{center}

% problems

\problem{1}

\subproblem{f}
\begin{gather*}
    M_5=20+30+32+26+12=120mC\\
    M_5=L_5=120mC
\end{gather*}
No, the agreement is an accident. $Q$ does lie inside the bracket.
\begin{gather*}
    E_L=\frac{5}{3}mC\\
    E_M=\frac{5}{3}mC\\
    E_R=\frac{25}{3}mC
\end{gather*}

\problem{3}

\subproblem{a}
\[
L_3=R_3=6+6+4=16L
\]
The agreement does not indicate that the result is accurate.\vem

\subproblem{b}
\begin{gather*}
    F'=-2t+3\\
    \max:(1.5,6.25)\\
    t_1 \implies \max : 6, \min :4\\
    t_2 \implies \max : 6.25, \min :6\\
    t_3 \implies \max : 6, \min :4\\
    \overline{S}_3=6+6.25+6=18.25L\\
    \underline{S}_3=4+6+4=14L  
\end{gather*}

\subproblem{c}
\begin{gather*}
    \overline{S}_3,\underline{S}_3 \implies [14,18.25], w=4.25L\\
    R_3,L_3\implies \{16\}, w=0L
\end{gather*}
The second approach assumes that the function is monotonically increasing over the interval.\vem

\subproblem{d}
\begin{gather*}
    16.5 \in [14,18.25]\\
    16.5 \notin \{16\}
\end{gather*}

\subproblem{e}
The second rectangle, as $F(t)$ changes monotonicity in the interval. \vem

\subproblem{f}
For functions that are monotonically increasing or decreasing. Otherwise, the lower and upper Riemann sums must be considered. 

\problem{4}

\subproblem{a}
\[
\overline{S}_n-\underline{S}_n=R_n-L_n=\sum_{k=1}^n f(x_{k})\Delta x-\sum_{k=1}^n f(x_{k-1})\Delta x=\sum_{k=1}^n \paren{f(x_k)-f(x_{k-1})}\Delta x
\]

\subproblem{b}
\begin{gather*}
    \overline{S}_4-\underline{S}_4=\Delta x\sum_{k=1}^4 \paren{f(x_k)-f(x_{k-1})}=\Delta x\paren{f(x_4)-f(x_0)}\\
    \overline{S}_n-\underline{S}_n=\sum_{k=1}^n \paren{f(x_k)-f(x_{k-1})}\Delta x=\paren{f(b)-f(a)}\Delta x=\frac{(b-a)\paren{f(b)-f(a)}}{n}
\end{gather*}

\subproblem{c}
\begin{gather*}
   \lim_{n \to \infty} \overline{S}_n-\underline{S}_n=\lim_{n \to \infty}\frac{(b-a)\paren{f(b)-f(a)}}{n} =0
\end{gather*}

\subproblem{d}
\begin{gather*}
 \overline{S}_n-\underline{S}_n=\frac{(b-a)\paren{f(b)-f(a)}}{n}\\
    \overline{S}_4-\underline{S}_4=\frac{(b-a)\paren{f(b)-f(a)}}{4}
\end{gather*}

\subproblem{e}
\begin{gather*}
    \overline{S}_n-\underline{S}_n=\frac{(b-a)\paren{f(b)-f(a)}}{n} < 0.001\\
    n>1000(b-a)\paren{f(b)-f(a)}\\
    \frac{\overline{S}_n-\underline{S}_n}{2}=\frac{(b-a)\paren{f(b)-f(a)}}{n} < 0.001\\
    n>500(b-a)\paren{f(b)-f(a)}
\end{gather*}
The existence of the identity reduces the amount of arithmetic needed to a minimal degree.\vem

\subproblem{f}
Identity for $f$ decreasing:
\[
\overline{S}_n-\underline{S}_n=\frac{(b-a)\paren{f(a)-f(b)}}{n}
\]
Problem 3's $F$ satisfies neither, as the function is not monotonic over the whole interval. The formulas and a proper amount of arithmetic are then needed to bound the value of the function.

\problem{5}

\[
D(x)=\begin{cases}
    1,\quad x \text{ rational}\\
    0,\quad x \text{ irrational}
\end{cases}
\]

\subproblem{a}
\begin{gather*}
    u\in \mathbb{Q} \implies \lim_{x \to 0+}u\in \mathbb{I}\\
    u\in \mathbb{I} \implies \lim_{x \to 0+}u\in \mathbb{Q}\\
    x\in [u,v] \implies D(x)\in \{0,1\}\\
    \text{sup }D=1, \text{inf } D=0\\
    \overline{S}_n=\sum_{k=1}^n1\cdot \frac{1}{n}=1\\
    \underline{S}_n=\sum_{k=1}^n0\cdot \frac{1}{n}=0\\
\end{gather*}

\subproblem{b}
\begin{gather*}
    S_{\mathbb{Q}}= \sum_{k=1}^n D\paren{\frac{k}{n}}= \sum_{k=1}^n \frac{1}{n}\cdot 1=1\\
    S_{\mathbb{I}}= \sum_{k=1}^n D\paren{\frac{k-1}{n}+\frac{1}{\sqrt{2}n}}= \sum_{k=1}^n \frac{1}{n}\cdot 0=0\\
\end{gather*}
The result of integrating $D(x)$ over $[0,1]$ solely depends on the set of points chosen, rendering no single value to be called as the result. \vem

\subproblem{c}
\begin{gather*}
    \overline{S}_n-\underline{S}_n=1-0=1
\end{gather*}
The difference between the sums remains as a constant no matter the size of $n$.
\[
\underline{S}_n<S_n<\overline{S}_n
\]
The range remains large and fails to be helpful towards more precise measurements. \vem

\subproblem{d}
\begin{gather*}
    T_n=\frac{1}{2}(L_n+R_n)=\frac{1}{2}(\overline{S}_n+\underline{S}_n)\\
    \underline{S}_n<A<\overline{S}_n\\
    |A-T_n| \le \frac{1}{2}(\overline{S}_n-\underline{S}_n)=\frac{(b-a)\paren{f(b)-f(a)}}{2n}
\end{gather*}

\subproblem{e}
\begin{gather*}
    \overline{S}_n-\underline{S}_n=\paren{f(b)-f(a)}\cdot \frac{(b-a)}{n}\le \frac{V(b-a)}{n}=V \Delta x\\
    f(x)=\sin x\implies V =1+1=2\\
    \overline{S}_n-\underline{S}_n\le \frac{2\pi}{n}\\
    \frac{2\pi}{n} \le \frac{1}{100} \implies n \ge 200\pi\\
    \boxed{n=629}
\end{gather*}





\end{document}

